\documentclass[10pt,a4paper]{amsart}
\usepackage{amsmath,amssymb,amsthm,graphicx}
\usepackage{verbatim}
\newenvironment{reference}{\comment}{\endcomment}
%\newenvironment{reference}{}{}
\newenvironment{slogan}{\comment}{\endcomment}
\newenvironment{history}{\comment}{\endcomment}

% For commutative diagrams we use Xy-pic
\usepackage[all]{xy}

% We use 2cell for 2-commutative diagrams.
\xyoption{2cell}
% Theorem environments.
%
\theoremstyle{plain}
\newtheorem{theorem}[subsection]{Theorem}
\newtheorem{proposition}[subsection]{Proposition}
\newtheorem{lemma}[subsection]{Lemma}

\theoremstyle{definition}
\newtheorem{definition}[subsection]{Definition}
\newtheorem{example}[subsection]{Example}
\newtheorem{exercise}[subsection]{Exercise}
\newtheorem{situation}[subsection]{Situation}

\theoremstyle{remark}
\newtheorem{remark}[subsection]{Remark}
\newtheorem{remarks}[subsection]{Remarks}

\numberwithin{equation}{subsection}

% Macros
%
\def\lim{\mathop{\mathrm{lim}}\nolimits}
\def\colim{\mathop{\mathrm{colim}}\nolimits}
\def\Spec{\mathop{\mathrm{Spec}}}
\def\Hom{\mathop{\mathrm{Hom}}\nolimits}
\def\Ext{\mathop{\mathrm{Ext}}\nolimits}
\def\SheafHom{\mathop{\mathcal{H}\!\mathit{om}}\nolimits}
\def\SheafExt{\mathop{\mathcal{E}\!\mathit{xt}}\nolimits}
\def\Sch{\mathit{Sch}}
\def\Mor{\mathop{\mathrm{Mor}}\nolimits}
\def\Ob{\mathop{\mathrm{Ob}}\nolimits}
\def\Sh{\mathop{\mathit{Sh}}\nolimits}
\def\NL{\mathop{N\!L}\nolimits}
\def\CH{\mathop{\mathrm{CH}}\nolimits}
\def\proetale{{pro\text{-}\acute{e}tale}}
\def\etale{{\acute{e}tale}}
\def\QCoh{\mathit{QCoh}}
\def\Ker{\mathop{\mathrm{Ker}}}
\def\Im{\mathop{\mathrm{Im}}}
\def\Coker{\mathop{\mathrm{Coker}}}
\def\Coim{\mathop{\mathrm{Coim}}}

% Boxtimes
%
\DeclareMathSymbol{\boxtimes}{\mathbin}{AMSa}{"02}

%
% Macros for moduli stacks/spaces
%
\def\QCohstack{\mathcal{QC}\!\mathit{oh}}
\def\Cohstack{\mathcal{C}\!\mathit{oh}}
\def\Spacesstack{\mathcal{S}\!\mathit{paces}}
\def\Quotfunctor{\mathrm{Quot}}
\def\Hilbfunctor{\mathrm{Hilb}}
\def\Curvesstack{\mathcal{C}\!\mathit{urves}}
\def\Polarizedstack{\mathcal{P}\!\mathit{olarized}}
\def\Complexesstack{\mathcal{C}\!\mathit{omplexes}}
% \Pic is the operator that assigns to X its picard group, usage \Pic(X)
% \Picardstack_{X/B} denotes the Picard stack of X over B
% \Picardfunctor_{X/B} denotes the Picard functor of X over B
\def\Pic{\mathop{\mathrm{Pic}}\nolimits}
\def\Picardstack{\mathcal{P}\!\mathit{ic}}
\def\Picardfunctor{\mathrm{Pic}}
\def\Deformationcategory{\mathcal{D}\!\mathit{ef}}
\usepackage{cleveref}
\crefname{originallemma}{original lemma}{original lemmas}
\Crefname{originallemma}{Original lemma}{Original lemmas}
\makeatletter
\newlabel{lemma-surjection-from-quasi-compact}{{040W}{source}{}{sourcepointer-1}{}}
\newlabel{lemma-surjection-from-quasi-compact@cref}{{[originallemma][0][]040W}{[1][0][]source}{}{}{}}
\makeatother
\crefname{originallemma}{original lemma}{original lemmas}
\Crefname{originallemma}{Original lemma}{Original lemmas}
\makeatletter
\newlabel{spaces-properties-lemma-space-locally-quasi-compact}{{04NN}{source}{}{sourcepointer-2}{}}
\newlabel{spaces-properties-lemma-space-locally-quasi-compact@cref}{{[originallemma][0][]04NN}{[1][0][]source}{}{}{}}
\makeatother
\title{Quasi-compactness descends along a universally open surjective morphism of algebraic spaces}
\author{The Stacks Project Authors}
\date{stacks-project:spaces-morphisms.tex:local-snapshot-20261012; GNU FDL 1.2 or later}
\begin{document}
\maketitle
\setcounter{subsection}{1}
\noindent\textit{Editorial note.} All delivered mathematical text is the author's own, copied from the pinned Stacks Project chapter file spaces-morphisms.tex; no statement or proof step was rewritten, reordered or removed.\par\medskip
\noindent\textit{Editorial note.} The proof is a descent argument on points and open covers; the local quasi-compactness of algebraic spaces is the only topological input, and the rest is the quasi-compactness criterion.\par\medskip
\noindent\textit{Editorial note.} The two references of the counted spans are recorded as non-load-bearing navigation pointers; the argument that consumes them stays inside the counted proof span.\par\medskip
\noindent\textit{Editorial note.} Preamble: the three native excerpt spans of sources/stacks-spaces-morphisms/preamble.tex (lines 8-13, 14-18 and 154-217). The chapter file itself carries no preamble text.\par\medskip
\noindent\textit{Allowed standard prerequisites (editorial).}\begin{itemize}
\item Quasi-compactness of a morphism of algebraic spaces tested after base change from quasi-compact spaces. Lemma lemma-characterize-representable-quasi-compact, used at proof lines 1360-1361 and 1372-1373. Is the criterion the proof verifies.
\item Points of an algebraic space have fundamental systems of quasi-compact open neighbourhoods. Properties of Spaces, Lemma spaces-properties-lemma-space-locally-quasi-compact, used at proof lines 1362-1367. Allows a quasi-compact open W surjecting onto Z to be found inside Y' x\_Y Z.
\item A surjective image of a quasi-compact space is quasi-compact. Lemma lemma-surjection-from-quasi-compact, used at proof lines 1370-1373. Transfers quasi-compactness from W x\_Y X to Z x\_Y X.
\item Openness and surjectivity are stable under base change. used at proof lines 1361-1362. Makes Y' x\_Y Z -> Z open and surjective.
\end{itemize}\medskip
\ifdefined\hypertarget\hypertarget{sourcepointer-1}{}\fi
\noindent\textit{Source reference.} Original lemma 040W: focused-seeds/AGP/sources/stacks-snapshot-20261011/tags.txt, source lines 5271--5271. Navigation pointer into the Stacks Project; the standard fact it supplies is applied inside the counted proof span and the argument resting on it stays there.\par\medskip
\ifdefined\hypertarget\hypertarget{sourcepointer-2}{}\fi
\noindent\textit{Source reference.} Original lemma 04NN: focused-seeds/AGP/sources/stacks-snapshot-20261011/tags.txt, source lines 6086--6086. Navigation pointer into the Stacks Project; the standard fact it supplies is applied inside the counted proof span and the argument resting on it stays there.\par\medskip
\begin{lemma}
\label{lemma-descent-quasi-compact}
Let $S$ be a scheme.
Let $f : X \to Y$ be a morphism of algebraic spaces over $S$.
Let $g : Y' \to Y$ be a universally open and surjective morphism of
algebraic spaces such that the base change $f' : X' \to Y'$ is quasi-compact.
Then $f$ is quasi-compact.
\end{lemma}

\begin{proof}
Let $Z \to Y$ be a morphism of algebraic spaces with $Z$ quasi-compact.
As $g$ is universally open and surjective, we see that
$Y' \times_Y Z \to Z$ is open and surjective. As every point of
$|Y' \times_Y Z|$ has a fundamental system of quasi-compact open
neighbourhoods (see
Properties of Spaces,
Lemma \ref{spaces-properties-lemma-space-locally-quasi-compact})
we can find a quasi-compact open $W \subset |Y' \times_Y Z|$
which surjects onto $Z$. Denote
$f'' : W \times_Y X \to W$ the base change of $f'$ by $W \to Y'$.
By assumption $W \times_Y X$ is quasi-compact. As $W \to Z$ is surjective
we see that $W \times_Y X \to Z \times_Y X$ is surjective.
Hence $Z \times_Y X$ is quasi-compact by
Lemma \ref{lemma-surjection-from-quasi-compact}.
Thus $f$ is quasi-compact.
\end{proof}

\end{document}
